\section{Conclusion}

Across four studies, language models use evidence selectively, infer
context-dependent response relationships, and benefit from training on related
relationships when forecasting unseen synthetic mechanisms and historical events.
These are the behavioral signatures of inductive forecasting defined in this work.

Forecasting ability is therefore not only a matter of matching outcomes. In these
studies, forecasts change systematically with evidence relevance, contextual
relationships, and training objectives. The gains are strongest when the relevant
relationship is recoverable from the task, and lower forecast error need not imply
recovery of the hidden mechanism.

The next test is whether these behaviors persist when cues are unreliable and the
relevant variables and mechanisms fall outside the training family.
